\section{Lean formalization and verification}\label{sec:lean}
Two standalone files build with Lean 4.19.0 and its standard library.
They give both a concrete field interface and a checked part of the
degree argument. The mathematical proof in the preceding sections
covers all even dimensions. The full object-theory derivation is a
remaining formalization task.

\paragraph{The two implemented layers.}
In \code{NamedField.lean}, a vector is
\code{Fin d} indexed with exact rational components.
The file defines the norms, distance, dot product, projection,
predicates (C), (U), (T), denominator $D$ (function \code{gridDen}),
the filtered integer grid, and the actual search
\code{firstWitness}. Its \code{GridHit} states that an accepted point
occurs on that grid. Its \code{TailBound} is the implication from an
accepted reading and (C) to the $25/32$ tail estimate.
The theorems \code{witness\_onBoundary} and \code{witness\_accepted}
extract verified properties of a returned point. The theorem
\code{assemble} proves the complete search statement conditional on
\code{GridHit} and \code{TailBound}. The file defines these two
obligations as propositions.

In \code{HairyBall.lean}, \code{finite\_obstruction} extracts a bounded
hit from unequal endpoint coefficients and preservation on the failure
branch. Its auxiliary index scan examines indices in descending order;
either enumeration order gives a good index. The theorem
\code{antipodal\_facet\_sign} proves the factor $(-1)^d$ by list
induction in arbitrary dimension. The namespace \code{Octahedron}
contains a more explicit verification for $S^2$. Its signed triangle
cochain \code{top} picks out the all-positive facet. For all $6^4$
ordered label quadruples the kernel verifies that its coboundary
vanishes when the quadruple contains no opposite pair.
\code{triangle\_prism} combines three such identities. Then
\code{chain\_prism} sums over any finite weighted triangular source,
and \code{degree\_step} gives equality of endpoint counts when the side
cochain vanishes on its boundary. \code{base\_cycle\_closed},
\code{base\_identity} and \code{base\_antipode} verify the boundary
and both counts for the eight original facets. These are ordinary Lean
theorems; the supplied files contain no admitted proofs.

\paragraph{Completing the argument in Lean.}
The following order puts a definite, testable statement at every interface.
\begin{enumerate}
\item \emph{Finite oriented chains.} Represent a complex by a duplicate-free
list of faces on ordered vertices, closed under subfaces. Implement
$\operatorname{orient}(u)$ as zero on duplicate tuples, else as the
parity of the inversion count times the increasing basis tuple.
Represent each degree's chain by a list of integer coefficients. Give
\code{boundary}, \code{push}, and \code{subdivide} the deletion,
pointwise-map and flag formulas of Section~\ref{sec:finite-degree}.
Prove $\partial f_\#=f_\#\partial$ by the two-deletion argument, and
$\partial S=S\partial$ by the explicit flag swap. A terminating
induction on the number of uncollected terms proves the signed-sum
pairing lemma. Use it for the boundary of $Z$ as well.
\item \emph{Degree obstruction.} Store the support map sending each fine
vertex to its original signed support. Prove the support and mesh bounds
when taking a rational barycentre. For a coarse $p$-face, prove
$c_\# S^r(s)=s$ by induction on its dimension, comparing a facet
coefficient. Implement the ordered prism formula with source tuples in
one global order. Prove the three deletion identities for the prism
operator, and that $d+1$ labels in a single cross-polytope face yield
zero. Applying these identities to $S^rZ$ gives the preservation
argument required by \code{finite\_obstruction}. In dimension two
\code{degree\_step} supplies the local-to-global step once the
subdivided source cycle is built.
\item \emph{Rational estimates.} Prove the triangle and reverse triangle
inequalities for \code{l1} by induction over the coordinate fold.
The identity $d\sum a_i^2-(\sum |a_i|)^2=\sum_{i<j}(|a_i|-|a_j|)^2$
follows by expansion or by adding the last coordinate. Use
cross multiplication with positive denominators to prove the projection
and normalized-edge bounds. Define the least-max label by folding over
indices, replacing only on strict improvement. The pair of inequalities
$1/(4d)<1/d$ and the signed-coordinate lower bound prove contiguity for
the two legs. In particular these proofs can run without a library of
real topology.
\item \emph{Completing the two named-field obligations.} Instantiate the
degree result on the $r$-fold subdivision in Section~\ref{sec:search}.
The projection gives exact tangency, and (T) controls its error.
Implement $\operatorname{round}_D$ by integer floor and prove its
distance bound by summing the rounding remainders. Prove membership in
\code{integerRows} by induction on the vector length, then apply the
filter lemma to the rounded vector. The displayed $3e/4$ bound now
proves \code{GridHit}. The coordinate triangle inequality and (C) give
\code{TailBound}. Apply \code{assemble}. For the round statement (R),
implement bisection by iterating the interval update and maintain
$l^2\le p\cdot p\le b^2$ and $b-l=2^{-j}$. The squared chord identity
and the integer chart formula yield the denominator and rounding bounds.
\item \emph{Embedding in the requested object theory.} Replace lists and
coefficient arrays by the pair/list codes in Section~\ref{sec:search},
rational comparisons by the corresponding PR Boolean terms, and each
fold by the equations of its fixed PR recursion. Use the repository's
\code{Core.Basic} induction and instantiation rules in
\code{IPRADerivations} \cite{deriv}. For example set
\[
 \mathsf{TabOK}(d,r,V,\tau)=
 \mathsf{MeshOK}(d,r)\ \mathbin{\wedge_B}\
 \mathsf{AllVerticesTangent}(d,r,V)\ \mathbin{\wedge_B}\
 \mathsf{AllEdgesSmall}(d,r,V,\tau).
 \]
Here the three factors respectively fold the rational tests
$\len{p-q}\le1/(8d)$ on simplex pairs, $p\cdot V[p]=0$ on
vertices, and $\len{V[p]-V[q]}\le\tau/(16d)$ on edges;
include type, shape and positivity guards at input. The closed numerical
endpoint asserts, for $d=2m+1$, that $\mathsf{TabOK}=1$ implies that
the scan index is in range and its norm test is true. Derive it by the
prefix invariants already given. For the main schema, specialize the
PR symbols for $A,\mu$, derive the finite table guard from the three
universally closed hypotheses, and instantiate the endpoint. Close the
resulting sentence in \code{Prf PRAAx []}. The existing bridge can
transport that derivation to Foundation's standard presentation. A
host-level Lean theorem about rational lists and the object derivation
are separate deliverables in this plan.
\end{enumerate}

ProveIt's \code{SphereBlueprint.lean} already defines signed vertices,
cross faces, the fundamental coefficient and its antipodal sign.
Its \code{SubdivisionNetTarget} and \code{BoundedApproximateZeroTarget}
are named obligations. It is possible to reuse the sign theorems when
porting the present proof to the repository toolchain \cite{spheres}.
The standalone files avoid a dependency on that checkout.

\Needspace{16\baselineskip}
\paragraph{A reproducible numerical example.}
The script \code{hairy\_ball.py} projects $b=(1,2,4)$ onto the plane
normal to the input direction. On the grid $D=16$ with requested
$\ell^1$ tolerance $2/5$, it returns
\[
 p=(-2,-5,-9)/16,\qquad V(p)=(7,-10,4)/55,\qquad
 \len{V(p)}=21/55<2/5.
\]
Here the dot product vanishes because $-14+50-36=0$. The exact polar
directions of this example have denominator $7$ in $P_\Q$, so the
positive residual at denominator $16$ is genuine. The degree demonstration
on that same grid uses $1{,}026$ vertices and $2{,}048$ oriented
triangles. With eight steps per leg the signed count drops from $1$
to $0$ between layers $6$ and $7$, and from $0$ to $-1$ between layers
$9$ and $10$. At the first drop the program exhibits a tetrahedron with
labels $(+2,+2,-2,-1)$. Thus one can see where the no-opposites condition
fails when normalizing near a zero on this coarse grid. The separate script
\code{round\_witness.py} implements rational stereographic rounding
and checks the distance to the true unit direction by exact comparisons
of quadratic radicals. Reproduction commands are in the archive README.
